% ============================================================
% THEO-CHIR-CHI-1: The Substrate Chirality Magnitude chi = phi^-3 as the
%   Locality-Selected Symmetric Bias of the 600-Cell's Two Nearest Shells
% A Layer-2/2.5 Structural-Selection Theorem (closes OPEN-CHIR-1d sub-gap 1d-alpha)
% Conscious Point Physics --- Substrate Chirality Arc
%   (series_umbrella/series_substrate_chirality_arc/chirality_derivations/)
%
% Purpose: justify the value of audit entry E21 (chi = phi^-3) by selecting the
%   distance-pair ratio that fixes its exponent. The substrate chirality magnitude
%   is the symmetric bias (d2-d1)/(d2+d1) of a 600-cell distance pair; the theorem
%   shows that the LOCALITY criterion -- the bias of the host vertex's two NEAREST
%   distance shells (the icosahedral vertex figure at d=1/phi and the dodecahedral
%   next shell at d=1, the local neighborhood where n-hat breaks H4 -> H3 = I_h) --
%   uniquely yields chi = phi^-3, and excludes the literature alternatives 1/sqrt5
%   and 5-2sqrt5 as NON-LOCAL (edge-to-distant-shell pairings). This closes the
%   ratio-selection sub-gap 1d-alpha of OPEN-CHIR-1d and answers the standing
%   reviewer question "why exponent -3?".
%
% Scope/layer: LAYER 2/2.5 structural-selection theorem. It does NOT eliminate
%   FI-C-9 (that requires the symmetry-breaking dynamics, sub-gap 1d-beta, deferred).
%   It selects chi = phi^-3 GIVEN (a) the symmetric-bias FORM of the magnitude
%   generator and (b) the locality criterion; the form itself is a stated structural
%   input. E21 stays emergent (provisional): 1d-alpha closed, 1d-beta open.
%
% Patch 0638 (Session 148) --- built from the Patch 0637 scope sketch
%   theo_chir_chi_magnitude_1_scope.md after the distance-spectrum exploration
%   (verify_chi_phi3_ratio.py) found the selection criterion. Audit-record
%   correction carried: E21 is derived (in value) by Capotauro Finding C-3, NOT by
%   THEO-CHIR-CONT-1.3 (which is magnitude inheritance, consumes chi).
%
% CHANGELOG
%   v1.0 (Patch 0638, Session 148, 29 May 2026) --- initial registration of
%     THEO-CHIR-CHI-1 at Layer 2/2.5, closing OPEN-CHIR-1d sub-gap 1d-alpha.
% ============================================================
% Version 1.2 --- 17 August 2026 (Patch 3213):
%   extended the generated identifier appendix to all five identifier
%   families (THEO-, FI-, PRED-, CONJ-, PH- as well as OPEN-), and
%   replaced review-convergence scores with AI review panel wording
%   where present. No claim, equation, number or section changed.
% Version 1.1 --- 17 August 2026 (Patch 3212):
%   added the generated plain-language appendix glossing the internal
%   programme identifiers used in this paper (founder jargon ruling, 16
%   Aug 2026). No claim, equation, number or section changed.

\documentclass[11pt,letterpaper]{article}
\usepackage[utf8]{inputenc}
\usepackage[margin=1in]{geometry}
\usepackage[T1]{fontenc}
\usepackage{amsmath,amssymb,amsfonts,amsthm}
\usepackage{xcolor}
\usepackage{hyperref}
\usepackage{enumitem}
\usepackage{parskip}
\usepackage{mdframed}
\usepackage{longtable}
\usepackage{array}

\hypersetup{colorlinks=true,linkcolor=blue,citecolor=blue,urlcolor=blue}

\newtheorem{theorem}{Theorem}[section]
\newtheorem{proposition}[theorem]{Proposition}
\newtheorem{lemma}[theorem]{Lemma}
\newtheorem{corollary}[theorem]{Corollary}
\newtheorem{definition}[theorem]{Definition}
\newtheorem{remark}[theorem]{Remark}

\newcommand{\nhat}{\hat{n}}
\newcommand{\vhost}{v_{\text{host}}}
\newcommand{\Hthree}{H_3}
\newcommand{\Hfour}{H_4}
\newcommand{\Icoh}{I_h}
\newcommand{\sbias}[2]{\frac{#2-#1}{#2+#1}}

\title{\textbf{THEO-CHIR-CHI-1}\\[4pt]
The Substrate Chirality Magnitude $\chi = \varphi^{-3}$ as the Locality-Selected
Symmetric Bias of the 600-Cell's Two Nearest Shells\\[6pt]
\large A Layer-2/2.5 Structural-Selection Theorem (OPEN-CHIR-1d, sub-gap 1d-$\alpha$)}
\author{Conscious Point Physics --- Substrate Chirality Arc}
\date{Version 1.0 --- Patch 0638, Session 148 --- 29 May 2026\\[2pt]
{\small Version~1.1 --- 17 August 2026 (Patch 3212: added the generated plain-language appendix glossing the internal programme identifiers used in this paper (founder jargon ruling, 16 Aug 2026). No claim, equation, number or section changed.)}\\[2pt]
{\small Version~1.2 --- 17 August 2026 (Patch 3213: extended the generated identifier appendix to all five identifier families (THEO-, FI-, PRED-, CONJ-, PH- as well as OPEN-), and replaced review-convergence scores with AI review panel wording where present. No claim, equation, number or section changed.)}}

\begin{document}
\maketitle

\begin{abstract}
\noindent
The substrate chirality magnitude $\chi = \varphi^{-3} \approx 0.236$ (audit entry E21)
is registered as a foundational input (FI-C-9) whose \emph{value} is computed, via
Capotauro Finding~C-3, as the symmetric bias $(1-\varphi^{-1})/(1+\varphi^{-1})$ of the
600-cell's edge-to-first-non-edge distance pair. The standing open question is not the
arithmetic but the \emph{ratio selection}: why this pair, rather than the literature
alternatives $1/\sqrt{5}$ and $5-2\sqrt{5}$? We enumerate the symmetric biases of all
$\binom{8}{2}=28$ distinct distance pairs in the 600-cell's distance spectrum from a host
vertex and show that a single \emph{locality} criterion --- the symmetric bias of the host
vertex's two \emph{nearest} distance shells (the $d=\varphi^{-1}$ icosahedral vertex figure
and the $d=1$ dodecahedral next shell, the local neighborhood where $\nhat$ breaks
$\Hfour \to \Hthree = \Icoh$) --- uniquely yields $\chi = \varphi^{-3}$. The exponent
$-3$ then follows as numerator $\varphi^{-2}$ over denominator $\varphi^{+1}$. The
alternatives are excluded as \emph{non-local}: $1/\sqrt{5}$ is the edge-to-$\varphi$-shell
bias (skipping three shells) and $5-2\sqrt{5}$ is the edge-to-antipode bias (the global
extreme); neither is an adjacent or two-nearest pairing. This closes the ratio-selection
sub-gap (1d-$\alpha$) of OPEN-CHIR-1d at Layer~2/2.5 and answers the reviewer's standing
``why exponent $-3$?'' question. It does \emph{not} eliminate FI-C-9: the value is selected
\emph{given} (a) the symmetric-bias \emph{form} of the magnitude generator and (b) the
locality criterion, and the symmetry-breaking \emph{dynamics} (sub-gap 1d-$\beta$) remain a
deferred deep target. E21 stays emergent (provisional): 1d-$\alpha$ closed, 1d-$\beta$ open.
All claims are verified at machine precision by \texttt{code/verify\_chi\_phi3\_ratio.py}.
\end{abstract}

\tableofcontents

% ============================================================
\section{Setup}
\label{sec:setup}
% ============================================================

Audit entry E21 is the substrate chirality magnitude $\chi = \varphi^{-3} \approx 0.236$,
controlling the substrate edge-length perturbation ($1 + \varepsilon\,\hat{e}\cdot\nhat$,
$|\varepsilon| = |\chi|$). Its value is computed by Capotauro Finding~C-3 as a
\emph{symmetric bias} of a 600-cell distance pair:
\begin{equation}
\chi \;=\; \sbias{\varphi^{-1}}{1} \;=\; \frac{\varphi^{-2}}{\varphi^{+1}} \;=\; \varphi^{-3},
\label{eq:findingC3}
\end{equation}
using $1 - \varphi^{-1} = \varphi^{-2}$ (since $\varphi^{-1} = \varphi-1 \Rightarrow
1-\varphi^{-1} = 2-\varphi = \varphi^{-2}$) and $1 + \varphi^{-1} = \varphi$. The
arithmetic is not in question. What is open --- and what a reviewer (ChatGPT, in the
THEO-CHIR-AUDIT-1 cycle) correctly pressed --- is the \emph{ratio selection}: the same
``symmetric bias of a distance pair'' template, applied to a different 600-cell pair, gives
$1/\sqrt{5} \approx 0.447$ or $5-2\sqrt{5} \approx 0.528$ (Capotauro \S6 Q3). Why
$\varphi^{-3}$?

\paragraph{Goal and non-goals.} This theorem closes the ratio-selection sub-gap (1d-$\alpha$)
of OPEN-CHIR-1d: it states a structural criterion that uniquely picks the
$\varphi^{-3}$-producing pair, and excludes the alternatives. It does \emph{not}: derive the
symmetry-breaking dynamics that would eliminate FI-C-9 (sub-gap 1d-$\beta$, deferred);
justify the symmetric-bias \emph{form} itself (a stated structural input, \S\ref{sec:form});
or treat the numerical signposts ($\delta_{CP}$, $\Delta p_{LR}$) as derivations
(\S\ref{sec:layer}). \emph{Record correction:} E21 is derived in value by Finding~C-3, not
by THEO-CHIR-CONT-1.3 (which is magnitude \emph{inheritance} --- it consumes $\chi$ through
the continuum-limit map and does not derive it).

% ============================================================
\section{The 600-cell distance spectrum at the host vertex}
\label{sec:spectrum}
% ============================================================

Take the 600-cell at unit circumradius (its 120 vertices the unit quaternions of the binary
icosahedral group $2I$). From a host vertex $\vhost$, the chord distance to any other vertex
is $\sqrt{2-2w}$, $w$ the scalar part relative to $\vhost$; the scalar parts take the nine
values $\{\pm1, \pm\varphi/2, \pm\tfrac12, \pm\varphi^{-1}/2, 0\}$, giving \emph{eight}
distinct non-zero distances. The spectrum (verified by
\texttt{verify\_chi\_phi3\_ratio.py}):

\begin{center}
\begin{tabular}{@{}rlrl@{}}
\hline
\textbf{shell} & \textbf{distance $d$} & \textbf{count} & \textbf{local-geometry role} \\
\hline
1 & $\varphi^{-1} \approx 0.618$ & 12 & icosahedron (the vertex figure) \\
2 & $1$ & 20 & dodecahedron (next shell; chord $=$ circumradius) \\
3 & $\sqrt{2-\varphi^{-1}} \approx 1.176$ & 12 & --- \\
4 & $\sqrt{2} \approx 1.414$ & 30 & equatorial \\
5 & $\varphi \approx 1.618$ & 12 & --- \\
6 & $\sqrt{3} \approx 1.732$ & 20 & --- \\
7 & $\sqrt{2+\varphi} \approx 1.902$ & 12 & --- \\
8 & $2$ & 1 & antipode \\
\hline
\end{tabular}
\end{center}

The two nearest shells are the $12$-vertex icosahedron (shell~1, the 600-cell's vertex
figure) and the $20$-vertex dodecahedron (shell~2). These are the dual pair forming the
\emph{local neighborhood} of $\vhost$: under vertex-aligned $\nhat = \vhost$, the substrate's
$\Hfour$ symmetry breaks to the residual $\Hthree = \Icoh$ at the host vertex (Capotauro
Reading~C), and shells 1--2 are precisely the icosahedral/dodecahedral structure carrying
that residual symmetry. Shell~2 sits at chord distance exactly $1$ (the circumradius).

% ============================================================
\section{The locality selection criterion}
\label{sec:selection}
% ============================================================

\begin{definition}[Symmetric-bias generator]
\label{def:bias}
For an ordered pair of distinct shell distances $d_i < d_j$, the \emph{symmetric bias} is
$\beta(d_i,d_j) = (d_j - d_i)/(d_j + d_i) \in (0,1)$. The chirality-magnitude generator is
taken to be a symmetric bias of a 600-cell distance pair (the structural \emph{form}; see
\S\ref{sec:form}).
\end{definition}

\begin{definition}[Locality criterion]
\label{def:locality}
A distance pair is \emph{local} if it consists of the host vertex's two nearest distance
shells (shells 1 and 2). Equivalently: the pair is \emph{adjacent} (consecutive shells)
\emph{and} is the shortest such pair (the two smallest distances).
\end{definition}

\begin{theorem}[Locality selects $\chi = \varphi^{-3}$; THEO-CHIR-CHI-1]
\label{thm:chi}
Among the symmetric biases of all $28$ distinct distance pairs of the 600-cell spectrum from
a host vertex (\S\ref{sec:spectrum}), the locality criterion (Definition~\ref{def:locality})
selects a \emph{unique} pair --- the two nearest shells $(\varphi^{-1}, 1)$ --- whose
symmetric bias is
\[
\beta(\varphi^{-1}, 1) \;=\; \sbias{\varphi^{-1}}{1} \;=\; \frac{\varphi^{-2}}{\varphi}
\;=\; \varphi^{-3}.
\]
Moreover $(\varphi^{-1},1)$ is the unique \emph{adjacent} pair whose bias equals
$\varphi^{-3}$, and it is the maximal-bias adjacent pair (bias $0.236$ vs.\ next-largest
adjacent bias $0.092$). The literature alternatives are \emph{non-local}: $1/\sqrt{5}$ is
produced only by the edge-to-$\varphi$-shell pair $(\varphi^{-1}, \varphi)$ (shells $1$--$5$,
non-adjacent), and $5-2\sqrt{5}$ only by the edge-to-antipode pair $(\varphi^{-1}, 2)$
(shells $1$--$8$, the global extreme). Hence the locality criterion fixes $\chi = \varphi^{-3}$
and excludes $1/\sqrt{5}$ and $5-2\sqrt{5}$.
\end{theorem}

\begin{proof}[Proof (enumeration).]
The eight distances of \S\ref{sec:spectrum} give $\binom{8}{2}=28$ pairs. Direct computation
of $\beta$ on all $28$ (script \texttt{verify\_chi\_phi3\_ratio.py}, machine precision) yields:
the seven \emph{adjacent} pairs have biases $\{0.236, 0.081, 0.092, 0.067, 0.034, 0.047,
0.025\}$ for shells $(1,2),(2,3),\dots,(7,8)$ respectively; the two-nearest pair $(1,2)$ is
the unique adjacent pair with bias $\varphi^{-3}=0.236$ and is the maximal adjacent bias.
Searching all $28$ pairs: $\beta = \varphi^{-3}$ occurs for $(\varphi^{-1},1)$ [adjacent,
two-nearest], $(1,\varphi)$ [shells $2$--$5$, non-adjacent], and
$(\sqrt{2-\varphi^{-1}},\sqrt{2+\varphi})$ [shells $3$--$7$, non-adjacent]; only the first is
local. $\beta = 1/\sqrt{5}$ occurs only for $(\varphi^{-1},\varphi)$ [shells $1$--$5$];
$\beta = 5-2\sqrt{5}$ occurs only for $(\varphi^{-1},2)$ [shells $1$--$8$]; both are
edge-to-distant-shell pairings, non-local. The locality criterion therefore selects
$(\varphi^{-1},1)$ uniquely, with bias $\varphi^{-3}$.
\end{proof}

\begin{remark}[The exponent]
The exponent $-3$ is fixed by the selected pair: numerator $1-\varphi^{-1} = \varphi^{-2}$,
denominator $1+\varphi^{-1} = \varphi^{+1}$, quotient $\varphi^{-3}$. ``Why $-3$'' reduces to
``why the two-nearest pair,'' answered by locality. The degeneracy of $\varphi^{-3}$ across
two further (non-local) pairs is a feature of the $\mathbb{Q}[\varphi]$ arithmetic of the
spectrum; locality picks the local representative.
\end{remark}

% ============================================================
\section{What is assumed: the symmetric-bias form}
\label{sec:form}
% ============================================================

Theorem~\ref{thm:chi} selects the ratio \emph{given} that the magnitude generator is a
symmetric bias $(d_j-d_i)/(d_j+d_i)$ of a distance pair (Definition~\ref{def:bias}). This
\emph{form} is a structural input, not derived here. Its motivation: a chirality magnitude is
a dimensionless left/right imbalance, and the symmetric bias is the natural scale-free
imbalance of two characteristic lengths; pairing the two \emph{nearest} lengths makes it a
\emph{local} substrate property at the symmetry-breaking locus. The honest status is that
1d-$\alpha$ resolves the \emph{ratio} selection within the symmetric-bias family; whether the
chirality magnitude must take the symmetric-bias form (vs.\ some other geometric functional of
the spectrum) is a residual structural assumption, recorded as a falsifier (\S\ref{sec:concl},
F2) rather than closed.

% ============================================================
\section{Layer, signposts, and what is not claimed}
\label{sec:layer}
% ============================================================

\paragraph{Layer 2/2.5.} This is a structural-selection argument over a finite combinatorial
object (the 600-cell distance spectrum), verified at machine precision. It is not a
dynamical derivation. It closes 1d-$\alpha$ (ratio selection) at Layer~2/2.5.

\paragraph{FI-C-9 not eliminated.} $\chi = \varphi^{-3}$ remains a foundational input until
the symmetry-breaking \emph{dynamics} (sub-gap 1d-$\beta$: which substrate-dynamical primitive
selects the broken chiral phase of $\Hfour \to I_4$) are derived. That is the deep deferred
target (F.1 \S14.17; deferred to the OPEN-SM-4 $\leftrightarrow$ SS-corpus). This theorem
strengthens E21 from ``value computed, ratio unjustified'' to ``value selected by a stated
locality criterion''; it does not retire FI-C-9.

\paragraph{Signposts are signposts.} The numerical proximities
$\delta_{CP} = 180^\circ + \arctan(\varphi^{-3}) \approx 193.3^\circ$ (vs.\ NuFIT
$195^\circ \pm 40^\circ$) and $\Delta p_{LR} = \varphi^{-3}/6 \approx 0.0394$ (vs.\ observed
$\sim 0.04$) \emph{support} $\varphi^{-3}$ but are not derivations (Capotauro Finding~C-5
discipline); they are not used as evidence in Theorem~\ref{thm:chi}.

\paragraph{E21 classification.} E21 stays \emph{emergent (provisional)}. 1d-$\alpha$ closed;
1d-$\beta$ open; the F.1 Phase-3 Case-A.1 ($\delta=\chi$) content that consumes $\chi$ remains
provisional. The classification does not change; its justification strengthens.

% ============================================================
\section{Conclusion and falsifiers}
\label{sec:concl}
% ============================================================

The substrate chirality magnitude $\chi = \varphi^{-3}$ is the locality-selected symmetric
bias of the 600-cell's two nearest distance shells (the icosahedral vertex figure at
$\varphi^{-1}$ and the dodecahedral shell at $1$), the local neighborhood where $\nhat$ breaks
$\Hfour \to \Hthree = \Icoh$. The exponent $-3$ follows; the alternatives $1/\sqrt{5}$ and
$5-2\sqrt{5}$ are excluded as non-local. Sub-gap 1d-$\alpha$ of OPEN-CHIR-1d is closed at
Layer~2/2.5; sub-gap 1d-$\beta$ (symmetry-breaking dynamics, FI-C-9 elimination) remains the
deferred deep target.

\paragraph{Falsifiers.} \textbf{(F1)} Exhibit a 600-cell distance pair satisfying the locality
criterion (Definition~\ref{def:locality}) with a symmetric bias $\neq \varphi^{-3}$ ---
impossible by the enumeration, so F1 is a check on the spectrum, not an open risk.
\textbf{(F2)} Show the chirality magnitude must take a \emph{non}-symmetric-bias form (some
other functional of the spectrum) yielding a value $\neq \varphi^{-3}$ --- this attacks the
\S\ref{sec:form} structural assumption, the genuine residual freedom in 1d-$\alpha$.
\textbf{(F3)} Derive the 1d-$\beta$ symmetry-breaking dynamics and obtain $\chi \neq \varphi^{-3}$
--- this would refute (not merely defer) the value, and is the decisive future test.

% ============================================================
\section*{References (internal)}
\addcontentsline{toc}{section}{References (internal)}
% ============================================================

\begin{itemize}[leftmargin=1.4em]
\item \texttt{chirality\_derivations/sketches/theo\_chir\_chi\_magnitude\_1\_scope.md} --- the
Patch 0637 scope sketch (1d-$\alpha$/1d-$\beta$ decomposition; the CONT-1.3$\to$Finding-C-3
correction).
\item \texttt{chirality\_derivations/code/verify\_chi\_phi3\_ratio.py} --- machine-precision
verification of the spectrum, the two-nearest shells, the $\varphi^{-3}$ bias, the adjacency
uniqueness, and the non-locality of the $1/\sqrt{5}$ and $5-2\sqrt{5}$ producers.
\item \texttt{capotauro/sketches/Capotauro\_chi\_phi\_closure.md} --- Finding~C-3 (the
symmetric-bias derivation of $\varphi^{-3}$; the $\varphi^{-2}$ archive-error correction);
Finding~C-4 (the eight distinct 600-cell distances); Findings C-5/C-7 ($\delta_{CP}$,
$\Delta p_{LR}$ signposts; $T = V/2 = 6$); \S6 Q3 (the candidate list $\varphi^{-3}$,
$1/\sqrt{5}$, $5-2\sqrt{5}$); FI-C-9.
\item \texttt{chirality\_audit/theo\_chir\_audit\_1.tex} (v1.1) --- E21 (emergent (P)); the
CONT-1.3 note corrected here.
\item \texttt{axiom-registry.md} --- A5 metric ($\eta = \ell_{\text{edge}}/R_{\text{circ}} =
1/\varphi$, the edge length entering the bias); A2 (600-cell distance structure).
\item \texttt{dynamical\_substrate\_law/sketches/F1\_subquestion\_pcd\_orientation\_link.md}
\S14.17 --- the $\chi$-from-pure-geometry long-term target ($\equiv$ 1d-$\beta$).
\item \texttt{frontier\_sectors/CHIR.md} --- OPEN-CHIR-1d; E21; THEO-CHIR-PCD-ORIENTATION-1
(the sibling E20 target).
\end{itemize}

\vfill
\noindent\rule{\textwidth}{0.4pt}\\
{\footnotesize THEO-CHIR-CHI-1 v1.0 --- Patch 0638, Session 148 --- Conscious Point Physics,
Substrate Chirality Arc. Layer-2/2.5 structural-selection theorem closing OPEN-CHIR-1d sub-gap
1d-$\alpha$: $\chi = \varphi^{-3}$ is the locality-selected symmetric bias of the 600-cell's
two nearest shells $(\varphi^{-1}, 1)$; alternatives $1/\sqrt{5}$, $5-2\sqrt{5}$ excluded as
non-local. Does not eliminate FI-C-9 (1d-$\beta$ dynamics deferred); E21 stays emergent
(provisional). Verified by \texttt{verify\_chi\_phi3\_ratio.py}.}

% BEGIN GENERATED IDENTIFIER APPENDIX -- do not edit by hand
\clearpage
\section*{Appendix: Programme identifiers used in this paper}
\addcontentsline{toc}{section}{Appendix: Programme identifiers}

\noindent This programme tracks its results, assumptions and unresolved questions under short identifiers, so that a claim and its dependencies can be cited precisely. Prefixes denote the kind of item: \texttt{THEO-} a proved theorem, \texttt{FI-} a foundational input the derivation assumes, \texttt{PRED-} a registered prediction, \texttt{CONJ-} a conjecture, \texttt{OPEN-} an unresolved question, \texttt{PH-} a problem history. Those appearing in this paper are glossed below; no external document is needed to read them.

\begin{description}
  \item[\texttt{FI-C-9}] \hfill \\ The substrate chirality magnitude, a free-magnitude postulate at v1.0 later grounded at v2.0.
  \item[\texttt{OPEN-CHIR-1}] \hfill \\ Derive each entry the audit classifies as *emergent* from
  \item[\texttt{OPEN-SM-4}] \hfill \\ Derive the lattice chirality-activation event that establishes χ = φ⁻³ and produces CP violation. *(φ⁻¹ was the original pre-Session-86 conjecture, superseded at Finding C-3 Patch 0378 — lost-1/φ arithmetic error φ⁻²→φ⁻³; the live magnitude is |χ| = φ⁻³ ≈ 0.236 per FI-C-9 / THEO-CHIR-CHI-1 / Capotauro v1.0–v2.0. Corrected Patch 0670.)*
  \item[\texttt{THEO-CHIR-AUDIT-1}] \hfill \\ The chirality structural audit theorem.
  \item[\texttt{THEO-CHIR-CHI-1}] \hfill \\ The substrate chirality magnitude result fixing the golden-ratio power of the chirality parameter.
  \item[\texttt{THEO-CHIR-CONT-1}] \hfill \\ Substrate-Handle-to-Effective-Coupling Bridge Theorem (first Layer 4 continuum-EFT projection closure under OPEN-SD-CHIR-PRIMITIVE umbrella; sub-claim (a) closure of joint Layer 4 paper §A shared bridge work at `series\_umbrella/series\_substrate\_chirality\_arc/chirality\_continuum/`; sector-agnostic continuum-EFT projection of the substrate handle \$\textbackslash\{\}chi/6\$ to chirality-sensitive effective operators in continuum EFT at observable scales)
  \item[\texttt{THEO-CHIR-PCD-ORIENTATION-1}] \hfill \\ A candidate theorem deriving the PCD winding orientation from the chirality director.
\end{description}
% END GENERATED IDENTIFIER APPENDIX

\end{document}
